% ============================================================
%  hayplot — Manual de Referência
%  Um plugin de visualização declarativa para a linguagem Hayashi
%  Compilar com: lualatex -interaction=nonstopmode manual.tex
% ============================================================
\documentclass[12pt, a4paper, oneside]{book}

% ── Idioma ─────────────────────────────────────────────────────────────────
\usepackage[brazilian]{babel}

% ── Tipografia (LuaLaTeX nativo) ────────────────────────────────────────────
\usepackage{fontspec}
\setmainfont{Latin Modern Roman}
\setmonofont{Latin Modern Mono}
\usepackage{microtype}

% ── Geometria e layout ──────────────────────────────────────────────────────
\usepackage[
  top=2.5cm, bottom=2.5cm,
  left=3cm,  right=2.5cm,
  headheight=15pt
]{geometry}
\usepackage{fancyhdr}
\pagestyle{fancy}
\fancyhf{}
\fancyhead[L]{\small\textit{\leftmark}}
\fancyhead[R]{\small\textit{hayplot}}
\fancyfoot[C]{\thepage}
\renewcommand{\headrulewidth}{0.4pt}

% ── Cores ───────────────────────────────────────────────────────────────────
\usepackage[dvipsnames, table]{xcolor}
\definecolor{hayblue}{HTML}{1565C0}
\definecolor{haygray}{HTML}{455A64}
\definecolor{haylight}{HTML}{E3F2FD}
\definecolor{codebg}{HTML}{F5F5F5}
\definecolor{haygreen}{HTML}{1B5E20}
\definecolor{hayred}{HTML}{B71C1C}
\definecolor{hayyellow}{HTML}{F57F17}

% ── Código-fonte ────────────────────────────────────────────────────────────
\usepackage{listings}
\lstdefinelanguage{hayashi}{
  morekeywords={let, const, fn, return, if, else, for, in, while, break,
                continue, import, load, generate, replace, display, print,
                export, write, dataframe, describe, predict, ols},
  sensitive=true,
  morestring=[b]",
  morecomment=[l]{//},
}
\lstset{
  language=hayashi,
  backgroundcolor=\color{codebg},
  basicstyle=\ttfamily\small,
  keywordstyle=\color{hayblue}\bfseries,
  stringstyle=\color{haygreen},
  commentstyle=\color{haygray}\itshape,
  numberstyle=\tiny\color{haygray},
  breaklines=true,
  breakatwhitespace=false,
  showstringspaces=false,
  tabsize=4,
  frame=single,
  framesep=4pt,
  framerule=0.4pt,
  rulecolor=\color{haygray!40},
  numbers=left,
  xleftmargin=1.5em,
  aboveskip=0.8em,
  belowskip=0.8em,
  captionpos=b,
}

% Estilo sem números para trechos inline
\lstdefinestyle{inline}{
  numbers=none,
  xleftmargin=0pt,
  frame=none,
  backgroundcolor=\color{codebg},
  basicstyle=\ttfamily\small,
}

% ── Hyperlinks ──────────────────────────────────────────────────────────────
\usepackage[
  colorlinks=true,
  linkcolor=hayblue,
  urlcolor=hayblue,
  citecolor=haygray,
  pdftitle={hayplot — Manual de Referência},
  pdfauthor={sheep-farm},
  pdfsubject={Plugin de visualização declarativa para Hayashi},
]{hyperref}

% ── Caixas informativas ─────────────────────────────────────────────────────
\usepackage{tcolorbox}
\tcbuselibrary{skins, breakable}

\newtcolorbox{infobox}[1][]{
  enhanced, breakable,
  colback=haylight, colframe=hayblue,
  fonttitle=\bfseries\small,
  left=6pt, right=6pt, top=4pt, bottom=4pt,
  #1
}
\newtcolorbox{warnbox}[1][]{
  enhanced, breakable,
  colback=hayyellow!10, colframe=hayyellow!70!black,
  fonttitle=\bfseries\small,
  left=6pt, right=6pt, top=4pt, bottom=4pt,
  #1
}
\newtcolorbox{tipbox}[1][]{
  enhanced, breakable,
  colback=green!6, colframe=haygreen,
  fonttitle=\bfseries\small,
  left=6pt, right=6pt, top=4pt, bottom=4pt,
  #1
}

% ── Tabelas ─────────────────────────────────────────────────────────────────
\usepackage{booktabs}
\usepackage{tabularx}
\usepackage{array}
\newcolumntype{L}[1]{>{\raggedright\arraybackslash}p{#1}}

% ── Ícones (sem fontawesome5) ──────────────────────────────────────────────
% Usa texto simples em vez de ícones para máxima portabilidade

% ── Índice ──────────────────────────────────────────────────────────────────
\usepackage{makeidx}
\makeindex
\usepackage{amssymb} % para \checkmark

% ── Atalhos ─────────────────────────────────────────────────────────────────
\newcommand{\code}[1]{\texttt{#1}}
\newcommand{\hayashi}{\textsc{Hayashi}}
\newcommand{\hayplot}{\textbf{hayplot}}

% ── Metadados da versão ─────────────────────────────────────────────────────
\newcommand{\hayversiontag}{v0.7.0}
\newcommand{\hayashiversiontag}{0.2.6}

% ============================================================
%  DOCUMENTO
% ============================================================
\begin{document}

% ── Página de título ────────────────────────────────────────────────────────
\begin{titlepage}
  \centering
  \vspace*{2.5cm}

  {\Huge\bfseries\color{hayblue} hayplot}\\[0.4em]
  {\large\color{haygray} Plugin de Visualização Declarativa para \hayashi}

  \vspace{1.5cm}

  \begin{tcolorbox}[
    enhanced, sharp corners,
    colback=hayblue, colframe=hayblue,
    width=0.85\textwidth,
    halign=center,
    fontupper=\color{white}\large
  ]
    \textbf{Grammar of Graphics} para a linguagem \hayashi \\[0.2em]
    \small Baseado em Apache Arrow FFI e Plotters (Rust puro)
  \end{tcolorbox}

  \vspace{2cm}

  \begin{tabularx}{0.65\textwidth}{@{}lX@{}}
    \toprule
    \textbf{Versão hayplot}     & \hayversiontag \\
    \textbf{Versão \hayashi{}}  & \hayashiversiontag \\
    \textbf{Repositório}        & \href{https://github.com/sheep-farm/hayplot}{\code{sheep-farm/hayplot}} \\
    \textbf{Licença}            & MIT \\
    \bottomrule
  \end{tabularx}

  \vfill
  {\small \today}
\end{titlepage}

% ── Sumário ─────────────────────────────────────────────────────────────────
\tableofcontents
\clearpage

% ============================================================
\chapter{Introdução}
% ============================================================

\hayplot{} é um plugin nativo para a linguagem \hayashi{} que implementa
uma \textbf{Gramática de Gráficos} (\textit{Grammar of Graphics}) ao estilo do
\code{ggplot2} do R. Ele permite construir visualizações de dados de forma
declarativa e composicional, usando o operador de pipe \code{|>} da linguagem.

\section{Filosofia}

A ideia central é que um gráfico pode ser descrito como a composição de
\emph{camadas semânticas independentes}:

\begin{enumerate}
  \item \textbf{Dados} — o DataFrame de entrada;
  \item \textbf{Mapeamento estético} (\textit{aes}) — quais colunas mapeiam para $x$ e $y$;
  \item \textbf{Geometrias} (\textit{geoms}) — como os dados serão representados visualmente;
  \item \textbf{Rótulos} (\textit{labs}) — título e nomes dos eixos.
\end{enumerate}

Cada camada é uma transformação pura de um dicionário de especificação, e a
renderização final acontece apenas ao chamar \code{render\_svg()}.

\section{Características Técnicas}

\begin{itemize}
  \item \textbf{Zero-Copy}: os dados do DataFrame são compartilhados com o plugin via
        Apache Arrow FFI sem cópia ou serialização --- o ponteiro Arrow é
        passado diretamente.

  \item \textbf{Rust puro}: usa a crate \code{plotters} compilada sem dependências
        externas de sistema (\textit{freetype}, \textit{fontconfig} não são necessários),
        garantindo portabilidade em Linux, macOS e Windows.

  \item \textbf{Saída SVG}: gera SVG vetorial escalável, adequado para publicação,
        relatórios em \LaTeX{} e inclusão em documentos web.

  \item \textbf{Composicional}: cada função recebe e retorna um \code{Dict}, tornando
        qualquer pipeline de plot testável, inspecionável e modificável.
\end{itemize}

\section{Instalação}
\index{instalação}

Instale o plugin diretamente do GitHub usando a CLI do \hayashi{}:

\begin{lstlisting}[style=inline, caption={Instalação via CLI}]
hay install sheep-farm/hayplot
\end{lstlisting}

O binário pré-compilado para a sua plataforma é baixado automaticamente
da última \textit{GitHub Release}. Para verificar se a instalação ocorreu com sucesso:

\begin{lstlisting}[style=inline, caption={Verificação da instalação}]
hay list
\end{lstlisting}

Você deve ver \code{sheep-farm/hayplot (native so plugin)} na lista de pacotes
instalados.

Para atualizar para a versão mais recente:

\begin{lstlisting}[style=inline, caption={Atualização}]
hay update sheep-farm/hayplot
\end{lstlisting}

Para verificar a integridade do que está instalado:

\begin{lstlisting}[style=inline, caption={Verificação de integridade}]
hay check-plugin sheep-farm/hayplot
\end{lstlisting}

Para remover o plugin:

\begin{lstlisting}[style=inline, caption={Remoção}]
hay remove sheep-farm/hayplot
\end{lstlisting}

\section{Importando o Plugin}
\index{import}

Em qualquer script \code{.hay}, importe o plugin com um alias:

\begin{lstlisting}[caption={Importação com alias}]
import("sheep-farm/hayplot", as=gg)
\end{lstlisting}

A partir daí, todas as funções são acessadas com o prefixo \code{gg::}.

% ============================================================
\chapter{Referência das Funções}
\index{funções}
% ============================================================

\hayplot{} expõe oito funções públicas. Cada uma recebe e retorna um
dicionário de especificação de plot (\textit{plot spec}), tornando a
composição via pipe natural e segura.

% ────────────────────────────────────────────────────────────
\section{\texttt{hayplot(df, aes)} --- Inicialização}
\index{hayplot()}
% ────────────────────────────────────────────────────────────

Inicializa um novo dicionário de especificação de plot com os dados e o
mapeamento estético.

\subsection*{Assinatura}

\begin{lstlisting}[style=inline]
hayplot(df: DataFrame, aes: Dict) -> Dict
\end{lstlisting}

\subsection*{Parâmetros}

\begin{tabularx}{\textwidth}{@{}lL{2.5cm}X@{}}
  \toprule
  \textbf{Parâmetro} & \textbf{Tipo} & \textbf{Descrição} \\
  \midrule
  \code{df}  & \code{DataFrame} & O DataFrame contendo as colunas a serem plotadas. \\[4pt]
  \code{aes} & \code{Dict}      & Mapeamento estético. Deve conter as chaves
                                  \code{"x"} e \code{"y"} com os nomes das colunas. \\
  \bottomrule
\end{tabularx}

\subsection*{Retorno}

Um \code{Dict} com a estrutura interna da especificação de plot:

\begin{lstlisting}[style=inline]
{
  "data":    <Arrow DataFrame>,
  "mapping": {"x": "coluna_x", "y": "coluna_y"},
  "layers":  [],
  "labs":    {}
}
\end{lstlisting}

\subsection*{Exemplo}

\begin{lstlisting}[caption={Criando a especificação inicial}]
let df = dataframe({"pib": [10, 20, 30], "vida": [70, 75, 80]})

let spec = gg::hayplot(df, {"x": "pib", "y": "vida"})
\end{lstlisting}

% ────────────────────────────────────────────────────────────
\section{\texttt{geom\_point(plot, color, size)} --- Camada de Pontos}
\index{geom\_point()}
% ────────────────────────────────────────────────────────────

Adiciona uma camada de dispersão (\textit{scatter plot}) à especificação.
Cada observação do DataFrame será representada por um círculo.

\subsection*{Assinatura}

\begin{lstlisting}[style=inline]
geom_point(plot: Dict, color: String, size: Float) -> Dict
\end{lstlisting}

\subsection*{Parâmetros}

\begin{tabularx}{\textwidth}{@{}lL{2.5cm}X@{}}
  \toprule
  \textbf{Parâmetro} & \textbf{Tipo} & \textbf{Descrição} \\
  \midrule
  \code{plot}  & \code{Dict}   & A especificação de plot retornada por \code{hayplot()}. \\[4pt]
  \code{color} & \code{String} & Cor dos pontos. Valores aceitos: \code{"red"}, \code{"blue"},
                                 \code{"green"}, \code{"yellow"}, \code{"magenta"},
                                 \code{"cyan"}, \code{"black"}. Padrão: \code{"blue"}. \\[4pt]
  \code{size}  & \code{Float}  & Raio dos pontos em pixels. Valores recomendados: 4.0--12.0. \\
  \bottomrule
\end{tabularx}

\subsection*{Retorno}

O mesmo \code{Dict} da especificação com a camada de pontos adicionada ao
campo \code{"layers"}.

\subsection*{Exemplo}

\begin{lstlisting}[caption={Adicionando uma camada de dispersão}]
let spec = gg::hayplot(df, {"x": "pib", "y": "vida"})
           |> gg::geom_point("red", 8.0)
\end{lstlisting}

\begin{warnbox}[title={\textbf{Atenção:} Nota sobre cores}]
  O nome da cor deve ser passado em inglês e minúsculas. Qualquer cor
  não reconhecida resultará na cor padrão \code{"blue"}.
  O suporte a cores hexadecimais (\code{\#RRGGBB}) está previsto para versões futuras.
\end{warnbox}

% ────────────────────────────────────────────────────────────
\section{\texttt{geom\_line(plot, color, size)} --- Camada de Linha}
\index{geom\_line()}
% ────────────────────────────────────────────────────────────

Adiciona uma camada de linha à especificação, conectando as observações do
DataFrame em ordem. Ideal para séries temporais, tendências e gráficos de evolução.
Pode ser combinado com \code{geom\_point} para criar gráficos de linhas com marcadores.

\subsection*{Assinatura}

\begin{lstlisting}[style=inline]
geom_line(plot: Dict, color: String, size: Float) -> Dict
\end{lstlisting}

\subsection*{Parâmetros}

\begin{tabularx}{\textwidth}{@{}lL{2.5cm}X@{}}
  \toprule
  \textbf{Parâmetro} & \textbf{Tipo} & \textbf{Descrição} \\
  \midrule
  \code{plot}  & \code{Dict}   & A especificação de plot retornada por \code{hayplot()}. \\[4pt]
  \code{color} & \code{String} & Cor da linha. Mesmos valores aceitos por \code{geom\_point}. \\[4pt]
  \code{size}  & \code{Float}  & Espessura da linha em pixels. Valores recomendados: 1.0--5.0. \\
  \bottomrule
\end{tabularx}

\subsection*{Retorno}

O mesmo \code{Dict} da especificação com a camada de linha adicionada ao
campo \code{"layers"}.

\subsection*{Composição de Camadas}

As camadas são renderizadas na \textbf{ordem em que são adicionadas}. Para
cobrir os pontos sobre a linha (evitando que a linha passe por cima),
adicione \code{geom\_line} \emph{antes} de \code{geom\_point}:

\begin{lstlisting}[caption={Gráfico de linha com marcadores de pontos}]
let plot = gg::hayplot(df, {"x": "mes", "y": "vendas"})
           |> gg::geom_line("blue", 3.0)    // linha primeiro
           |> gg::geom_point("red", 6.0)   // pontos por cima
           |> gg::labs("Vendas Mensais", "Mes", "Vendas")
\end{lstlisting}

% ────────────────────────────────────────────────────────────
\section{\texttt{geom\_bar(plot, color, width)} --- Camada de Barras}
\index{geom\_bar()}
% ────────────────────────────────────────────────────────────

Adiciona uma camada de gráfico de barras à especificação. Cada observação
do DataFrame será representada por uma barra vertical posicionada no valor $x$
com altura correspondente ao valor $y$.

\subsection*{Assinatura}

\begin{lstlisting}[style=inline]
geom_bar(plot: Dict, color: String, width: Float) -> Dict
\end{lstlisting}

\subsection*{Parâmetros}

\begin{tabularx}{\textwidth}{@{}lL{2.5cm}X@{}}
  \toprule
  \textbf{Parâmetro} & \textbf{Tipo} & \textbf{Descrição} \\
  \midrule
  \code{plot}  & \code{Dict}   & A especificação de plot retornada por \code{hayplot()}. \\[4pt]
  \code{color} & \code{String} & Cor das barras. Mesmos valores aceitos por \code{geom\_point}. \\[4pt]
  \code{width} & \code{Float}  & Largura das barras centradas no valor $x$. Valores recomendados: 0.4--1.0. \\
  \bottomrule
\end{tabularx}

\subsection*{Retorno}

O mesmo \code{Dict} da especificação com a camada de barras adicionada ao
campo \code{"layers"}.

\subsection*{Exemplo}

\begin{lstlisting}[caption={Gráfico de barras categórico}]
let plot = gg::hayplot(df, {"x": "categoria", "y": "vendas"})
           |> gg::geom_bar("blue", 0.6)
           |> gg::labs("Vendas por Categoria", "Categoria", "Vendas")
\end{lstlisting}

% ────────────────────────────────────────────────────────────
\section{\texttt{geom\_histogram(plot, color, bins)} --- Camada de Histograma}
\index{geom\_histogram()}
% ────────────────────────────────────────────────────────────

Adiciona uma camada de histograma à especificação. Calcula automaticamente
a distribuição de frequência dos valores da coluna $y$ e divide em um número
especificado de \textit{bins} (intervalos).

\subsection*{Assinatura}

\begin{lstlisting}[style=inline]
geom_histogram(plot: Dict, color: String, bins: Int) -> Dict
\end{lstlisting}

\subsection*{Parâmetros}

\begin{tabularx}{\textwidth}{@{}lL{2.5cm}X@{}}
  \toprule
  \textbf{Parâmetro} & \textbf{Tipo} & \textbf{Descrição} \\
  \midrule
  \code{plot}  & \code{Dict}   & A especificação de plot retornada por \code{hayplot()}. \\[4pt]
  \code{color} & \code{String} & Cor das barras do histograma. Mesmos valores aceitos por \code{geom\_point}. \\[4pt]
  \code{bins}  & \code{Int}    & Número de intervalos (\textit{bins}) para dividir os dados. Valores recomendados: 5--30. \\
  \bottomrule
\end{tabularx}

\subsection*{Retorno}

O mesmo \code{Dict} da especificação com a camada de histograma adicionada ao
campo \code{"layers"}.

\subsection*{Exemplo}

\begin{lstlisting}[caption={Histograma de distribuição}]
let plot = gg::hayplot(df, {"x": "dummy", "y": "notas"})
           |> gg::geom_histogram("green", 15)
           |> gg::labs("Distribuição de Notas", "Nota", "Frequência")
\end{lstlisting}

\begin{tipbox}[title={\textbf{Dica:} Dummy $x$ para histogramas}]
  Para histogramas univariados, use uma coluna $x$ constante (e.g., \code{"dummy": [1.0, 1.0, ...]})
  pois o histograma calcula a distribuição apenas a partir dos valores $y$.
\end{tipbox}

% ────────────────────────────────────────────────────────────
\section{\texttt{geom\_boxplot(plot, color, width)} --- Camada de Boxplot}
\index{geom\_boxplot()}
% ────────────────────────────────────────────────────────────

Adiciona uma camada de boxplot (\textit{box-and-whisker plot}) à especificação.
Exibe estatísticas resumidas: quartis (Q1, Q3), mediana, e \textit{whiskers}
(1.5$\times$IQR). Ideal para visualizar distribuição e identificar outliers.

\subsection*{Assinatura}

\begin{lstlisting}[style=inline]
geom_boxplot(plot: Dict, color: String, width: Float) -> Dict
\end{lstlisting}

\subsection*{Parâmetros}

\begin{tabularx}{\textwidth}{@{}lL{2.5cm}X@{}}
  \toprule
  \textbf{Parâmetro} & \textbf{Tipo} & \textbf{Descrição} \\
  \midrule
  \code{plot}  & \code{Dict}   & A especificação de plot retornada por \code{hayplot()}. \\[4pt]
  \code{color} & \code{String} & Cor do box e whiskers. Mesmos valores aceitos por \code{geom\_point}. \\[4pt]
  \code{width} & \code{Float}  & Largura da caixa. Valores recomendados: 0.3--0.7. \\
  \bottomrule
\end{tabularx}

\subsection*{Retorno}

O mesmo \code{Dict} da especificação com a camada de boxplot adicionada ao
campo \code{"layers"}.

\subsection*{Estatísticas Calculadas}

\begin{itemize}
  \item \textbf{Q1}: 25º percentil (limite inferior da caixa)
  \item \textbf{Mediana}: 50º percentil (linha preta na caixa)
  \item \textbf{Q3}: 75º percentil (limite superior da caixa)
  \item \textbf{IQR}: Q3 $-$ Q1 (intervalo interquartil)
  \item \textbf{Whiskers}: Q1 $-$ 1.5$\times$IQR e Q3 $+$ 1.5$\times$IQR (limites dos whiskers)
\end{itemize}

\subsection*{Exemplo}

\begin{lstlisting}[caption={Boxplot de distribuição salarial}]
let plot = gg::hayplot(df, {"x": "departamento", "y": "salario"})
           |> gg::geom_boxplot("red", 0.4)
           |> gg::labs("Distribuição Salarial", "Departamento", "Salário (USD)")
\end{lstlisting}

% ────────────────────────────────────────────────────────────
\section{\texttt{geom\_heatmap(plot, color, cell\_size)} --- Camada de Heatmap}
\index{geom\_heatmap()}
% ────────────────────────────────────────────────────────────

Adiciona uma camada de heatmap (\textit{mapa de calor}) à especificação.
Cada par $(x, y)$ é renderizado como uma célula retangular cuja intensidade de cor
é proporcional ao valor $y$ normalizado. Ideal para visualizar dados em grade 2D.

\subsection*{Assinatura}

\begin{lstlisting}[style=inline]
geom_heatmap(plot: Dict, color: String, cell_size: Float) -> Dict
\end{lstlisting}

\subsection*{Parâmetros}

\begin{tabularx}{\textwidth}{@{}lL{2.5cm}X@{}}
  \toprule
  \textbf{Parâmetro} & \textbf{Tipo} & \textbf{Descrição} \\
  \midrule
  \code{plot}      & \code{Dict}   & A especificação de plot retornada por \code{hayplot()}. \\[4pt]
  \code{color}     & \code{String} & Cor base para o heatmap. A intensidade é calculada misturando com branco. \\[4pt]
  \code{cell\_size} & \code{Float}  & Tamanho das células do heatmap. Valores recomendados: 0.5--1.5. \\
  \bottomrule
\end{tabularx}

\subsection*{Retorno}

O mesmo \code{Dict} da especificação com a camada de heatmap adicionada ao
campo \code{"layers"}.

\subsection*{Normalização de Cores}

Os valores $y$ são normalizados para o intervalo $[0, 1]$ usando a fórmula:

\[
\text{intensidade} = \frac{y - y_{\min}}{y_{\max} - y_{\min}}
\]

A cor final é calculada misturando a cor base com branco proporcionalmente à
intensidade (maior intensidade = cor mais saturada).

\subsection*{Exemplo}

\begin{lstlisting}[caption={Heatmap de temperatura em grade}]
let plot = gg::hayplot(df, {"x": "coord_x", "y": "coord_y"})
           |> gg::geom_heatmap("red", 0.8)
           |> gg::labs("Temperatura", "Coordenada X", "Coordenada Y")
\end{lstlisting}

% ────────────────────────────────────────────────────────────
\section{\texttt{geom\_area(plot, color, size)} --- Camada de Área}
\index{geom\_area()}
% ────────────────────────────────────────────────────────────

Adiciona uma camada de gráfico de área à especificação. Preenche a área sob a linha
conectando os pontos, ordenando automaticamente por $x$. Ideal para visualizar valores
acumulados, volume ao longo do tempo ou PnL.

\subsection*{Assinatura}

\begin{lstlisting}[style=inline]
geom_area(plot: Dict, color: String, size: Float) -> Dict
\end{lstlisting}

\subsection*{Parâmetros}

\begin{tabularx}{\textwidth}{@{}lL{2.5cm}X@{}}
  \toprule
  \textbf{Parâmetro} & \textbf{Tipo} & \textbf{Descrição} \\
  \midrule
  \code{plot}  & \code{Dict}   & A especificação de plot retornada por \code{hayplot()}. \\[4pt]
  \code{color} & \code{String} & Cor da área e da linha. Mesmos valores aceitos por \code{geom\_point}. \\[4pt]
  \code{size}  & \code{Float}  & Espessura da linha de contorno. Valores recomendados: 1.0--3.0. \\
  \bottomrule
\end{tabularx}

\subsection*{Retorno}

O mesmo \code{Dict} da especificação com a camada de área adicionada ao
campo \code{"layers"}.

\subsection*{Comportamento}

Os pontos são automaticamente ordenados por $x$ antes do rendering. A área é
preenchida usando uma aproximação por polígonos entre pontos consecutivos, com
base em $y=0$ como limite inferior.

\subsection*{Exemplo}

\begin{lstlisting}[caption={Gráfico de área para receita acumulada}]
let plot = gg::hayplot(df, {"x": "mes", "y": "receita_acumulada"})
           |> gg::geom_area("blue", 2.0)
           |> gg::labs("Receita Acumulada", "Mês", "Receita (milhares)")
\end{lstlisting}

\begin{tipbox}[title={\textbf{Dica:} Ordenação automática}]
  Diferente de \code{geom\_line}, \code{geom\_area} ordena automaticamente os pontos
  por $x$ antes de renderizar, garantindo que a área seja preenchida corretamente
  mesmo se os dados não estiverem ordenados.
\end{tipbox}

% ────────────────────────────────────────────────────────────
\section{\texttt{geom\_hline(plot, color, size, yintercept)} --- Linha Horizontal}
\index{geom\_hline()}
% ────────────────────────────────────────────────────────────

Adiciona uma linha horizontal de referência ao gráfico em um valor $y$ específico.
Ideal para marcar thresholds, limites de risco, médias ou valores alvo.

\subsection*{Assinatura}

\begin{lstlisting}[style=inline]
geom_hline(plot: Dict, color: String, size: Float, yintercept: Float) -> Dict
\end{lstlisting}

\subsection*{Parâmetros}

\begin{tabularx}{\textwidth}{@{}lL{2.5cm}X@{}}
  \toprule
  \textbf{Parâmetro} & \textbf{Tipo} & \textbf{Descrição} \\
  \midrule
  \code{plot}       & \code{Dict}   & A especificação de plot retornada por \code{hayplot()}. \\[4pt]
  \code{color}      & \code{String} & Cor da linha. Mesmos valores aceitos por \code{geom\_point}. \\[4pt]
  \code{size}       & \code{Float}  & Espessura da linha em pixels. Valores recomendados: 0.5--2.0. \\[4pt]
  \code{yintercept} & \code{Float}  & Valor $y$ onde a linha horizontal será desenhada. \\
  \bottomrule
\end{tabularx}

\subsection*{Exemplo}

\begin{lstlisting}[caption={Linha horizontal de referência}]
let plot = gg::hayplot(df, {"x": "tempo", "y": "valor"})
           |> gg::geom_line("blue", 2.0)
           |> gg::geom_hline("red", 1.0, 50.0)  // Linha em y=50
           |> gg::labs("Métrica com Threshold", "Tempo", "Valor")
\end{lstlisting}

% ────────────────────────────────────────────────────────────
\section{\texttt{geom\_vline(plot, color, size, xintercept)} --- Linha Vertical}
\index{geom\_vline()}
% ────────────────────────────────────────────────────────────

Adiciona uma linha vertical de referência ao gráfico em um valor $x$ específico.
Ideal para marcar eventos, split de ações, marcos temporais ou pontos de corte.

\subsection*{Assinatura}

\begin{lstlisting}[style=inline]
geom_vline(plot: Dict, color: String, size: Float, xintercept: Float) -> Dict
\end{lstlisting}

\subsection*{Parâmetros}

\begin{tabularx}{\textwidth}{@{}lL{2.5cm}X@{}}
  \toprule
  \textbf{Parâmetro} & \textbf{Tipo} & \textbf{Descrição} \\
  \midrule
  \code{plot}       & \code{Dict}   & A especificação de plot retornada por \code{hayplot()}. \\[4pt]
  \code{color}      & \code{String} & Cor da linha. Mesmos valores aceitos por \code{geom\_point}. \\[4pt]
  \code{size}       & \code{Float}  & Espessura da linha em pixels. Valores recomendados: 0.5--2.0. \\[4pt]
  \code{xintercept} & \code{Float}  & Valor $x$ onde a linha vertical será desenhada. \\
  \bottomrule
\end{tabularx}

\subsection*{Exemplo}

\begin{lstlisting}[caption={Linha vertical de evento}]
let plot = gg::hayplot(df, {"x": "ano", "y": "receita"})
           |> gg::geom_line("blue", 2.0)
           |> gg::geom_vline("green", 1.0, 2020.0)  // Linha em x=2020
           |> gg::labs("Receita com Evento Marcado", "Ano", "Receita")
\end{lstlisting}

% ────────────────────────────────────────────────────────────
\section{\texttt{geom\_abline(plot, color, size, slope, intercept)} --- Linha Diagonal}
\index{geom\_abline()}
% ────────────────────────────────────────────────────────────

Adiciona uma linha diagonal de referência seguindo a equação $y = \text{slope} \times x + \text{intercept}$.
Ideal para linhas de 45° (benchmark), tendências teóricas ou relações lineares esperadas.

\subsection*{Assinatura}

\begin{lstlisting}[style=inline]
geom_abline(plot: Dict, color: String, size: Float, slope: Float, intercept: Float) -> Dict
\end{lstlisting}

\subsection*{Parâmetros}

\begin{tabularx}{\textwidth}{@{}lL{2.5cm}X@{}}
  \toprule
  \textbf{Parâmetro} & \textbf{Tipo} & \textbf{Descrição} \\
  \midrule
  \code{plot}      & \code{Dict}   & A especificação de plot retornada por \code{hayplot()}. \\[4pt]
  \code{color}     & \code{String} & Cor da linha. Mesmos valores aceitos por \code{geom\_point}. \\[4pt]
  \code{size}      & \code{Float}  & Espessura da linha em pixels. Valores recomendados: 0.5--2.0. \\[4pt]
  \code{slope}     & \code{Float}  & Coeficiente angular da linha. \\[4pt]
  \code{intercept} & \code{Float}  & Intercepto da linha (valor em $y$ quando $x=0$). \\
  \bottomrule
\end{tabularx}

\subsection*{Exemplo}

\begin{lstlisting}[caption={Linha diagonal de benchmark]}
let plot = gg::hayplot(df, {"x": "x", "y": "y"})
           |> gg::geom_point("blue", 5.0)
           |> gg::geom_abline("red", 1.0, 1.0, 0.0)  // Linha y = x
           |> gg::labs("Observado vs. Esperado", "X", "Y")
\end{lstlisting}

% ────────────────────────────────────────────────────────────
\section{\texttt{geom\_step(plot, color, size, direction)} --- Linha em Degrau}
\index{geom\_step()}
% ────────────────────────────────────────────────────────────

Adiciona uma linha em degrau conectando os pontos. Útil para funções de distribuição acumulada,
preços discretos ou dados que mudam em passos. A direção controla se o degrau é horizontal-então-vertical
(\texttt{"hv"}) ou vertical-então-horizontal (\texttt{"vh"}).

\subsection*{Assinatura}

\begin{lstlisting}[style=inline]
geom_step(plot: Dict, color: String, size: Float, direction: String) -> Dict
\end{lstlisting}

\subsection*{Parâmetros}

\begin{tabularx}{\textwidth}{@{}lL{2.5cm}X@{}}
  \toprule
  \textbf{Parâmetro} & \textbf{Tipo} & \textbf{Descrição} \\
  \midrule
  \code{plot}      & \code{Dict}   & A especificação de plot retornada por \code{hayplot()}. \\[4pt]
  \code{color}     & \code{String} & Cor da linha. Mesmos valores aceitos por \code{geom\_point}. \\[4pt]
  \code{size}      & \code{Float}  & Espessura da linha em pixels. Valores recomendados: 1.0--3.0. \\[4pt]
  \code{direction} & \code{String} & Direção do degrau: \code{"hv"} (horizontal-então-vertical) ou \code{"vh"} (vertical-então-horizontal). \\
  \bottomrule
\end{tabularx}

\subsection*{Exemplo}

\begin{lstlisting}[caption={Gráfico de degrau para preços}]
let plot = gg::hayplot(df, {"x": "tempo", "y": "preco"})
           |> gg::geom_step("blue", 2.0, "hv")
           |> gg::labs("Mudanças de Preço (Degrau)", "Tempo", "Preço")
\end{lstlisting}

\section{\texttt{scale\_x\_log10(plot)} --- Escala Logarítmica em X}
\index{scale\_x\_log10()}
% ───────────────────────────────────────────────────────────---

Define o eixo $x$ para escala logarítmica base 10. Ideal para dados com crescimento exponencial
ou ranges muito amplos. Valores não-positivos ($\le 0$) são tratados como \texttt{NaN}.

\subsection*{Assinatura}

\begin{lstlisting}[style=inline]
scale_x_log10(plot: Dict) -> Dict
\end{lstlisting}

\subsection*{Parâmetros}

\begin{tabularx}{\textwidth}{@{}lL{2.5cm}X@{}}
  \toprule
  \textbf{Parâmetro} & \textbf{Tipo} & \textbf{Descrição} \\
  \midrule
  \code{plot} & \code{Dict} & A especificação de plot retornada por \code{hayplot()}. \\
  \bottomrule
\end{tabularx}

\subsection*{Exemplo}

\begin{lstlisting}[caption={Eixo X em escala logarítmica}]
let plot = gg::hayplot(df, {"x": "frequencia", "y": "amplitude"})
           |> gg::geom_point("blue", 5.0)
           |> gg::scale_x_log10()
           |> gg::labs("Resposta em Frequência", "Frequência (log10)", "Amplitude")
\end{lstlisting}

% ───────────────────────────────────────────────────────────---
\section{\texttt{scale\_y\_log10(plot)} --- Escala Logarítmica em Y}
\index{scale\_y\_log10()}
% ───────────────────────────────────────────────────────────---

Define o eixo $y$ para escala logarítmica base 10. Ideal para dados com crescimento exponencial
ou ranges muito amplos. Valores não-positivos ($\le 0$) são tratados como \texttt{NaN}.

\subsection*{Assinatura}

\begin{lstlisting}[style=inline]
scale_y_log10(plot: Dict) -> Dict
\end{lstlisting}

\subsection*{Parâmetros}

\begin{tabularx}{\textwidth}{@{}lL{2.5cm}X@{}}
  \toprule
  \textbf{Parâmetro} & \textbf{Tipo} & \textbf{Descrição} \\
  \midrule
  \code{plot} & \code{Dict} & A especificação de plot retornada por \code{hayplot()}. \\
  \bottomrule
\end{tabularx}

\subsection*{Exemplo}

\begin{lstlisting}[caption={Eixo Y em escala logarítmica}]
let plot = gg::hayplot(df, {"x": "tempo", "y": "valor"})
           |> gg::geom_line("red", 2.0)
           |> gg::scale_y_log10()
           |> gg::labs("Crescimento Exponencial", "Tempo", "Valor (log10)")
\end{lstlisting}

\begin{tipbox}[title={\textbf{Dica:} Cores hexadecimais}]
  Todos os parâmetros de cor aceitam tanto cores nomeadas (\code{"red"}, \code{"blue"}, etc.)
  quanto códigos hexadecimais no formato \code{\#RRGGBB} (ex: \code{"\#FF5733"}, \code{"\#C70039"}).
  Isso permite especificar cores personalizadas com precisão.
\end{tipbox}

% ─────────────────────────────────────────────────────────---
\section{\texttt{facet\_wrap(plot, group\_col)} --- Faceting por Grupo}
\index{facet\_wrap()}
% ─────────────────────────────────────────────────────────---

Especifica uma coluna para faceting. Quando usado com \code{render\_facets()}, gera
múltiplos gráficos SVG separados, um para cada valor único na coluna de grupo.
Esta é uma abordagem simplificada de faceting — painéis combinados em um único SVG
não são suportados; cada grupo gera seu próprio arquivo SVG.

\subsection*{Assinatura}

\begin{lstlisting}[style=inline]
facet_wrap(plot: Dict, group_col: String) -> Dict
\end{lstlisting}

\subsection*{Parâmetros}

\begin{tabularx}{\textwidth}{@{}lL{2.5cm}X@{}}
  \toprule
  \textbf{Parâmetro} & \textbf{Tipo} & \textbf{Descrição} \\
  \midrule
  \code{plot}      & \code{Dict}   & A especificação de plot retornada por \code{hayplot()}. \\[4pt]
  \code{group\_col} & \code{String} & Nome da coluna numérica usada para agrupar os dados. \\
  \bottomrule
\end{tabularx}

\subsection*{Exemplo}

\begin{lstlisting}[caption={Faceting por grupo}]
let plot = gg::hayplot(df, {"x": "tempo", "y": "valor"})
           |> gg::facet_wrap("categoria")
           |> gg::geom_point("blue", 5.0)
           |> gg::labs("Dados por Categoria", "Tempo", "Valor")
\end{lstlisting}

% ─────────────────────────────────────────────────────────---
\section{\texttt{render\_facets(plot)} --- Renderizar Painéis}
\index{render\_facets()}
% ─────────────────────────────────────────────────────────---

Renderiza múltiplos gráficos SVG separados, um para cada grupo especificado por \code{facet\_wrap()}.
Retorna uma lista de strings SVG. O usuário deve salvar cada SVG em arquivos separados.

\subsection*{Assinatura}

\begin{lstlisting}[style=inline]
render_facets(plot: Dict) -> Result<List<String>, String>
\end{lstlisting}

\subsection*{Parâmetros}

\begin{tabularx}{\textwidth}{@{}lL{2.5cm}X@{}}
  \toprule
  \textbf{Parâmetro} & \textbf{Tipo} & \textbf{Descrição} \\
  \midrule
  \code{plot} & \code{Dict} & A especificação de plot com \code{facet\_wrap()} aplicado. \\
  \bottomrule
\end{tabularx}

\subsection*{Retorno}

Uma lista (\code{List<String>}) contendo as strings SVG para cada grupo, ordenadas pelo valor do grupo.

\subsection*{Exemplo}

\begin{lstlisting}[caption={Renderizar e salvar painéis faceted}]
let plot = gg::hayplot(df, {"x": "tempo", "y": "valor"})
           |> gg::facet_wrap("categoria")
           |> gg::geom_point("blue", 5.0)
           |> gg::labs("Dados por Categoria", "Tempo", "Valor")

let svg_list = gg::render_facets(plot)
// svg_list[0] contém SVG para grupo 1
// svg_list[1] contém SVG para grupo 2
// ...
\end{lstlisting}

\begin{tipbox}[title={\textbf{Nota:} Faceting simplificado}]
  Esta implementação gera arquivos SVG separados em vez de painéis combinados.
  Para visualizar todos os grupos simultaneamente, abra os SVGs gerados em tabs do navegador
  ou use ferramentas externas para combinar as imagens.
\end{tipbox}

% ─────────────────────────────────────────────────────────---

% ────────────────────────────────────────────────────────────
\section{\texttt{labs(plot, title, x, y)} --- Rótulos}
\index{labs()}
% ────────────────────────────────────────────────────────────

Define o título do gráfico e os rótulos dos eixos horizontal e vertical.

\subsection*{Assinatura}

\begin{lstlisting}[style=inline]
labs(plot: Dict, title: String, x: String, y: String) -> Dict
\end{lstlisting}

\subsection*{Parâmetros}

\begin{tabularx}{\textwidth}{@{}lL{2.5cm}X@{}}
  \toprule
  \textbf{Parâmetro} & \textbf{Tipo} & \textbf{Descrição} \\
  \midrule
  \code{plot}  & \code{Dict}   & A especificação de plot. \\[4pt]
  \code{title} & \code{String} & Título exibido no topo do gráfico. \\[4pt]
  \code{x}     & \code{String} & Rótulo do eixo horizontal. \\[4pt]
  \code{y}     & \code{String} & Rótulo do eixo vertical. \\
  \bottomrule
\end{tabularx}

\subsection*{Retorno}

O mesmo \code{Dict} com o campo \code{"labs"} preenchido.

\subsection*{Exemplo}

\begin{lstlisting}[caption={Definindo título e nomes dos eixos}]
let spec = gg::hayplot(df, {"x": "pib", "y": "vida"})
           |> gg::geom_point("blue", 6.0)
           |> gg::labs("PIB vs. Expectativa de Vida",
                       "PIB per Capita (USD)",
                       "Expectativa de Vida (anos)")
\end{lstlisting}

% ────────────────────────────────────────────────────────────
\section{\texttt{render\_svg(plot)} --- Renderização}
\index{render\_svg()}
% ────────────────────────────────────────────────────────────

Compila a especificação de plot e renderiza o gráfico final como uma
\textbf{string SVG}. Esta é a única função que acessa os dados do DataFrame
e executa operações de desenho.

\subsection*{Assinatura}

\begin{lstlisting}[style=inline]
render_svg(plot: Dict) -> Result<String, String>
\end{lstlisting}

\subsection*{Parâmetros}

\begin{tabularx}{\textwidth}{@{}lL{2.5cm}X@{}}
  \toprule
  \textbf{Parâmetro} & \textbf{Tipo} & \textbf{Descrição} \\
  \midrule
  \code{plot} & \code{Dict} & A especificação de plot completa (com dados, mapeamento,
                              camadas e rótulos). \\
  \bottomrule
\end{tabularx}

\subsection*{Retorno}

\begin{itemize}
  \item Em caso de sucesso: uma \code{String} contendo o código XML do SVG
        (800$\times$600 pixels).
  \item Em caso de erro: uma mensagem de erro descritiva como \code{String}.
\end{itemize}

\begin{infobox}[title={\textbf{Nota:} Dimensões do SVG}]
  A saída SVG tem resolução fixa de \textbf{800$\times$600 pixels}. Como é um
  formato vetorial, pode ser redimensionada sem perda de qualidade em qualquer
  aplicação ou editor de imagens.
\end{infobox}

\subsection*{Processo interno de renderização}

Internamente, \code{render\_svg()} executa os seguintes passos:

\begin{enumerate}
  \item Extrai o ponteiro Arrow do campo \code{"data"} via FFI Zero-Copy;
  \item Lê o mapeamento \code{"x"} e \code{"y"} e recupera as colunas do DataFrame;
  \item Calcula os limites automáticos dos eixos com 10\% de margem;
  \item Itera sobre \code{"layers"} e aplica cada geometria;
  \item Compila o gráfico para uma string SVG em memória usando \code{plotters}.
\end{enumerate}

\subsection*{Exemplo}

\begin{lstlisting}[caption={Renderizando e salvando o SVG}]
let svg_content = gg::render_svg(plot)
write(svg_content, "saida.svg")
\end{lstlisting}

% ============================================================
\chapter{Fluxo de Trabalho}
% ============================================================

\section{Pipeline Declarativo}

O fluxo típico de um gráfico em \hayplot{} segue a estrutura abaixo:

\begin{lstlisting}[caption={Estrutura geral de um pipeline hayplot}]
import("sheep-farm/hayplot", as=gg)

// 1. Carregar ou criar os dados
load "dados.csv" as df

// 2. Inicializar a especificacao com dados e mapeamento
// 3. Adicionar camadas geometricas
// 4. Definir rotulos
let plot = gg::hayplot(df, {"x": "col_x", "y": "col_y"})
           |> gg::geom_point("blue", 6.0)
           |> gg::labs("Titulo do Grafico", "Eixo X", "Eixo Y")

// 5. Renderizar e salvar
let svg = gg::render_svg(plot)
write(svg, "grafico.svg")
\end{lstlisting}

\section{Tipos de Colunas Suportados}

A função \code{render\_svg()} aceita colunas numéricas nos seguintes tipos Arrow:

\begin{center}
  \begin{tabular}{@{}ll@{}}
    \toprule
    \textbf{Tipo Arrow} & \textbf{Tipo Hayashi} \\
    \midrule
    \code{Float64}      & \code{float} \\
    \code{Int64}        & \code{int} \\
    \bottomrule
  \end{tabular}
\end{center}

\begin{warnbox}[title={\textbf{Atenção:} Colunas categóricas}]
  Colunas do tipo \code{String} ou booleanas \textbf{não} são suportadas como
  coordenadas $x$ ou $y$. Para variáveis categóricas, encode-as numericamente
  antes de plotar.
\end{warnbox}

\section{Valores Ausentes}

Observações com \code{NaN} em qualquer uma das coordenadas são automaticamente
omitidas da renderização. Não é necessário filtrar os dados antes de plotar.

% ============================================================
\chapter{Exemplos Completos}
% ============================================================

\section{Dispersão Simples}
\index{exemplo!dispersão simples}

O exemplo mais básico: dados criados inline, sem arquivos externos.

\begin{lstlisting}[caption={examples/basic\_scatter.hay}]
import("sheep-farm/hayplot", as=gg)

// 1. Criar dataset PIB vs Expectativa de Vida
let data = {
    "gdp_per_capita":  [12000, 24000, 35000, 48000, 65000, 85000],
    "life_expectancy": [68.5,  72.1,  76.4,  79.2,  81.5,  83.1]
}
let df = dataframe(data)

// 2. Construir a especificacao do plot
print("Building basic scatter plot...")
let plot = gg::hayplot(df, {"x": "gdp_per_capita", "y": "life_expectancy"})
           |> gg::geom_point("blue", 7.0)
           |> gg::labs(
               "GDP vs Life Expectancy",
               "GDP per Capita (USD)",
               "Life Expectancy (Years)"
           )

// 3. Renderizar e salvar
let svg = gg::render_svg(plot)
write(svg, "gdp_life_exp.svg")
print("Plot saved to 'gdp_life_exp.svg'!")
\end{lstlisting}

\section{Curva de Phillips}
\index{exemplo!curva de Phillips}

Clássico da macroeconomia: relação negativa entre desemprego e inflação.

\begin{lstlisting}[caption={examples/phillips\_curve.hay}]
import("sheep-farm/hayplot", as=gg)

// 1. Dataset macroeconomico simulado
let data = {
    "unemployment_rate": [3.5, 4.2, 5.0, 6.1, 7.5, 8.8, 9.5],
    "inflation_rate":    [5.2, 4.5, 3.8, 2.9, 1.8, 1.2, 0.9]
}
let df = dataframe(data)

// 2. Construir o grafico da Curva de Phillips
print("Building Phillips Curve plot...")
let plot = gg::hayplot(df, {"x": "unemployment_rate", "y": "inflation_rate"})
           |> gg::geom_point("red", 8.0)
           |> gg::labs(
               "The Phillips Curve",
               "Unemployment Rate (%)",
               "Inflation Rate (%)"
           )

// 3. Renderizar e salvar
let svg = gg::render_svg(plot)
write(svg, "phillips_curve.svg")
print("Plot saved to 'phillips_curve.svg'!")
\end{lstlisting}

\section{A partir de Arquivo CSV}
\index{exemplo!CSV}

Integrando \hayplot{} com o comando \code{load} do \hayashi{}:

\begin{lstlisting}[caption={Plot a partir de um arquivo externo}]
import("sheep-farm/hayplot", as=gg)

// 1. Carregar dados externos
load "dados_painel.csv" as df

// 2. Criar grafico
let plot = gg::hayplot(df, {"x": "renda", "y": "consumo"})
           |> gg::geom_point("green", 5.0)
           |> gg::labs(
               "Renda vs. Consumo",
               "Renda (R$)",
               "Consumo (R$)"
           )

// 3. Renderizar
let svg = gg::render_svg(plot)
write(svg, "renda_consumo.svg")
\end{lstlisting}

\section{Combinando com Regressão OLS}
\index{exemplo!OLS}

\hayplot{} integra-se naturalmente ao workflow econométrico do \hayashi{}.
O exemplo a seguir estima um modelo OLS e visualiza os valores ajustados
contra os observados:

\begin{lstlisting}[caption={Visualizando valores ajustados de uma regressão OLS}]
import("sheep-farm/hayplot", as=gg)

// 1. Carregar dados e estimar OLS
load "painel.csv" as df
let m = ols(Y ~ X1 + X2, df)

// 2. Gerar valores preditos
predict df yhat = m

// 3. Scatter: observado vs predito
let plot = gg::hayplot(df, {"x": "yhat", "y": "Y"})
           |> gg::geom_point("blue", 5.0)
           |> gg::labs(
               "Observado vs. Ajustado",
               "Y-hat (valores preditos)",
               "Y (valores observados)"
           )

// 4. Salvar
let svg = gg::render_svg(plot)
write(svg, "obs_vs_fit.svg")
\end{lstlisting}

\section{Pipeline Multi-linha}
\index{exemplo!pipeline multi-linha}

Com o suporte a quebra de linha no operador \code{|>} do \hayashi{}, é
possível escrever pipelines mais legíveis:

\begin{lstlisting}[caption={Usando quebras de linha na pipeline}]
import("sheep-farm/hayplot", as=gg)

let df = dataframe({"x": [1.0, 2.0, 3.0, 4.0, 5.0],
                    "y": [2.5, 4.1, 5.9, 8.2, 10.0]})

let svg = gg::hayplot(df, {"x": "x", "y": "y"})
    |> gg::geom_point("cyan", 9.0)
    |> gg::labs("Tendencia Linear", "Variavel X", "Variavel Y")
    |> gg::render_svg

write(svg, "tendencia.svg")
\end{lstlisting}

\section{Linha com Marcadores de Pontos}
\index{exemplo!linha com pontos}

Combinando \code{geom\_line} e \code{geom\_point} em uma única pipeline, é
possível criar gráficos de tendência com marcadores visuais em cada observação.
A \textbf{ordem das camadas} determina qual é desenhada por cima:

\begin{lstlisting}[caption={examples/line\_and\_scatter.hay}]
import("sheep-farm/hayplot", as=gg)

// Dados de crescimento de vendas mensais
let data = {"month": [1.0, 2.0, 3.0, 4.0, 5.0, 6.0, 7.0, 8.0, 9.0, 10.0, 11.0, 12.0],
            "sales": [10.5, 12.0, 11.2, 14.8, 16.5, 15.0, 18.2, 21.0, 19.5, 22.8, 25.0, 24.2]}
let df = dataframe(data)

// geom_line adicionado antes de geom_point para que os pontos fiquem por cima
let plot = gg::hayplot(df, {"x": "month", "y": "sales"})
    |> gg::geom_line("blue", 3.0)
    |> gg::geom_point("red", 6.0)
    |> gg::labs("Sales Growth Over Time", "Month", "Sales (Thousands)")

let svg_content = gg::render_svg(plot)
write(svg_content, "sales_growth.svg")
print("Plot saved to 'sales_growth.svg'!")
\end{lstlisting}

\section{Gráfico de Barras}
\index{exemplo!barras}

Visualização de dados categóricos ou agregados por grupos:

\begin{lstlisting}[caption={examples/bar\_chart.hay}]
import("sheep-farm/hayplot", as=gg)

// Dados de vendas por categoria
let data = {"category": [1.0, 2.0, 3.0, 4.0, 5.0],
            "sales": [150.0, 230.0, 180.0, 320.0, 290.0]}
let df = dataframe(data)

let plot = gg::hayplot(df, {"x": "category", "y": "sales"})
    |> gg::geom_bar("blue", 0.6)
    |> gg::labs("Sales by Category", "Category", "Sales (Thousands)")

let svg_content = gg::render_svg(plot)
write(svg_content, "bar_chart.svg")
print("Plot saved to 'bar_chart.svg'!")
\end{lstlisting}

\section{Histograma}
\index{exemplo!histograma}

Análise de distribuição de uma variável univariada:

\begin{lstlisting}[caption={examples/histogram.hay}]
import("sheep-farm/hayplot", as=gg)

// Dados de notas de teste
let data = {"dummy": [1.0, 1.0, 1.0, 1.0, 1.0],
            "scores": [65.0, 72.0, 78.0, 85.0, 88.0, 90.0, 92.0, 95.0,
                       78.0, 82.0, 75.0, 88.0, 91.0, 84.0, 79.0, 86.0,
                       93.0, 77.0, 83.0, 89.0]}
let df = dataframe(data)

let plot = gg::hayplot(df, {"x": "dummy", "y": "scores"})
    |> gg::geom_histogram("green", 15)
    |> gg::labs("Distribution of Test Scores", "Score", "Frequency")

let svg_content = gg::render_svg(plot)
write(svg_content, "histogram.svg")
print("Plot saved to 'histogram.svg'!")
\end{lstlisting}

\section{Boxplot}
\index{exemplo!boxplot}

Visualização de distribuição e identificação de outliers:

\begin{lstlisting}[caption={examples/boxplot.hay}]
import("sheep-farm/hayplot", as=gg)

// Dados salariais por departamento
let data = {"department": [1.0, 1.0, 1.0, 1.0, 1.0],
            "salary": [45000.0, 52000.0, 48000.0, 61000.0, 55000.0,
                       58000.0, 49000.0, 63000.0, 51000.0, 57000.0]}
let df = dataframe(data)

let plot = gg::hayplot(df, {"x": "department", "y": "salary"})
    |> gg::geom_boxplot("red", 0.4)
    |> gg::labs("Salary Distribution", "Department", "Salary (USD)")

let svg_content = gg::render_svg(plot)
write(svg_content, "boxplot.svg")
print("Plot saved to 'boxplot.svg'!")
\end{lstlisting}

\section{Heatmap}
\index{exemplo!heatmap}

Visualização de dados em grade 2D com intensidade de cor:

\begin{lstlisting}[caption={examples/heatmap.hay}]
import("sheep-farm/hayplot", as=gg)

// Dados de temperatura em grade 3x3
let data = {"x": [1.0, 2.0, 3.0, 1.0, 2.0, 3.0, 1.0, 2.0, 3.0],
            "y": [1.0, 1.0, 1.0, 2.0, 2.0, 2.0, 3.0, 3.0, 3.0],
            "temperature": [20.0, 25.0, 30.0, 22.0, 28.0, 35.0,
                            18.0, 24.0, 32.0]}
let df = dataframe(data)

let plot = gg::hayplot(df, {"x": "x", "y": "y"})
    |> gg::geom_heatmap("red", 0.8)
    |> gg::labs("Temperature Heatmap", "X Location", "Y Location")

let svg_content = gg::render_svg(plot)
write(svg_content, "heatmap.svg")
print("Plot saved to 'heatmap.svg'!")
\end{lstlisting}

\section{Gráfico de Área}
\index{exemplo!área}

Visualização de valores acumulados com área preenchida:

\begin{lstlisting}[caption={examples/area\_chart.hay}]
import("sheep-farm/hayplot", as=gg)

// Dados de receita acumulada mensal
let data = {"month": [1.0, 2.0, 3.0, 4.0, 5.0, 6.0, 7.0, 8.0, 9.0, 10.0, 11.0, 12.0],
            "revenue": [10.0, 25.0, 45.0, 70.0, 100.0, 135.0, 175.0, 220.0,
                        270.0, 325.0, 385.0, 450.0]}
let df = dataframe(data)

let plot = gg::hayplot(df, {"x": "month", "y": "revenue"})
    |> gg::geom_area("blue", 2.0)
    |> gg::labs("Cumulative Revenue Over Time", "Month", "Revenue (Thousands)")

let svg_content = gg::render_svg(plot)
write(svg_content, "area_chart.svg")
print("Plot saved to 'area_chart.svg'!")
\end{lstlisting}

\section{Linhas de Referência}
\index{exemplo!linhas de referência}

Combinação de linhas horizontais, verticais e diagonais para anotações:

\begin{lstlisting}[caption={examples/reference\_lines.hay}]
import("sheep-farm/hayplot", as=gg)

// Dados de vendas mensais
let data = {"month": [1.0, 2.0, 3.0, 4.0, 5.0, 6.0, 7.0, 8.0, 9.0, 10.0, 11.0, 12.0],
            "sales": [10.5, 12.0, 11.2, 14.8, 16.5, 15.0, 18.2, 21.0,
                        19.5, 22.8, 25.0, 24.2]}
let df = dataframe(data)

let plot = gg::hayplot(df, {"x": "month", "y": "sales"})
    |> gg::geom_line("blue", 2.0)
    |> gg::geom_hline("red", 1.0, 15.0)    // Horizontal line at y=15
    |> gg::geom_vline("green", 1.0, 6.0)   // Vertical line at x=6
    |> gg::geom_abline("magenta", 1.0, 1.5, 8.0)  // Diagonal line: y = 1.5x + 8
    |> gg::labs("Sales with Reference Lines", "Month", "Sales (Thousands)")

let svg_content = gg::render_svg(plot)
write(svg_content, "reference_lines.svg")
print("Plot saved to 'reference_lines.svg'!")
\end{lstlisting}

\section{Gráfico de Degrau}
\index{exemplo!degrau}

Visualização de mudanças discretas com linhas em degrau:

\begin{lstlisting}[caption={examples/step\_chart.hay}]
import("sheep-farm/hayplot", as=gg)

// Dados de preços em tempo discreto
let data = {"time": [1.0, 2.0, 3.0, 4.0, 5.0, 6.0, 7.0, 8.0],
            "price": [100.0, 105.0, 103.0, 108.0, 106.0, 110.0, 109.0, 112.0]}
let df = dataframe(data)

let plot = gg::hayplot(df, {"x": "time", "y": "price"})
    |> gg::geom_step("blue", 2.0, "hv")  // horizontal then vertical
    |> gg::labs("Price Changes (Step Chart)", "Time", "Price")

let svg_content = gg::render_svg(plot)
write(svg_content, "step_chart.svg")
print("Plot saved to 'step_chart.svg'!")
\end{lstlisting}

\section{Escala Logarítmica}
\index{exemplo!escala logarítmica}

Visualização de dados exponenciais com escala logarítmica e cores hexadecimais:

\begin{lstlisting}[caption={examples/log\_scale.hay}]
import("sheep-farm/hayplot", as=gg)

// Dados com crescimento exponencial
let data = {"x": [1.0, 10.0, 100.0, 1000.0, 10000.0, 100000.0],
            "y": [1.0, 100.0, 10000.0, 1000000.0, 100000000.0, 10000000000.0]}
let df = dataframe(data)

let plot = gg::hayplot(df, {"x": "x", "y": "y"})
    |> gg::geom_point("#FF5733", 5.0)  // Custom hex color
    |> gg::geom_line("#C70039", 2.0)
    |> gg::scale_x_log10()
    |> gg::scale_y_log10()
    |> gg::labs("Exponential Growth (Log Scale)", "X (log10)", "Y (log10)")

let svg_content = gg::render_svg(plot)
write(svg_content, "log_scale.svg")
print("Plot saved to 'log_scale.svg'!")
\end{lstlisting}

\section{Faceting por Grupo}
\index{exemplo!faceting}

Visualização de dados agrupados com gráficos separados:

\begin{lstlisting}[caption={examples/facet\_wrap.hay}]
import("sheep-farm/hayplot", as=gg)

// Dados com grupos
let data = {"x": [1.0, 2.0, 3.0, 1.0, 2.0, 3.0, 1.0, 2.0, 3.0],
            "y": [10.0, 20.0, 30.0, 15.0, 25.0, 35.0, 12.0, 22.0, 32.0],
            "group": [1.0, 1.0, 1.0, 2.0, 2.0, 2.0, 3.0, 3.0, 3.0]}
let df = dataframe(data)

let plot = gg::hayplot(df, {"x": "x", "y": "y"})
    |> gg::facet_wrap("group")
    |> gg::geom_point("blue", 5.0)
    |> gg::geom_line("red", 2.0)
    |> gg::labs("Grouped Data", "X", "Y")

let svg_list = gg::render_facets(plot)
// svg_list contains SVG strings for each group
print("Generated " + svg_list.length + " facet plots")
\end{lstlisting}

% ============================================================
\chapter{Arquitetura Técnica}
% ============================================================

\section{Comunicação via FFI Arrow}

\hayplot{} é um \textbf{plugin nativo} do \hayashi{}: uma biblioteca dinâmica
(\code{.so}/\code{.dll}/\code{.dylib}) carregada em tempo de execução.
A comunicação entre o interpretador e o plugin segue o protocolo de
FFI do \code{hayashi-plugin-sdk}:

\begin{enumerate}
  \item O interpretador serializa os argumentos em uma estrutura C-ABI;
  \item DataFrames são passados como ponteiros Arrow \code{ArrayRef}, sem cópia dos dados;
  \item O plugin processa em Rust nativo e retorna os valores;
  \item O interpretador desserializa o retorno e libera a memória via hooks de desalocação.
\end{enumerate}

\begin{infobox}[title={\textbf{Nota:} Por que Zero-Copy?}]
  DataFrames econométricos podem ter milhões de linhas. Passar os dados
  via serialização (JSON, por exemplo) causaria duplicação de memória e
  \textit{overhead} significativo de CPU. O Arrow FFI resolve isso com um
  protocolo de ponteiros C padronizado pela Apache Software Foundation.
\end{infobox}

\section{Dependências}

\begin{center}
  \begin{tabular}{@{}llL{5cm}@{}}
    \toprule
    \textbf{Crate}             & \textbf{Versão} & \textbf{Função} \\
    \midrule
    \code{hayashi-plugin-sdk}  & \textit{main}   & FFI ABI, tipos Arrow, macros de exportação \\
    \code{plotters}            & \code{0.3.7}    & Renderização SVG, primitivas gráficas \\
    \code{arrow}               & (via SDK)       & Tipos de dados colunares \\
    \bottomrule
  \end{tabular}
\end{center}

\code{plotters} é compilado com \code{default-features = false} e apenas com
as features \code{svg\_backend}, \code{area\_series} e \code{line\_series},
eliminando todas as dependências de sistema.

\section{Compilação e Distribuição}

O plugin é compilado como \code{cdylib} (biblioteca C dinâmica):

\begin{lstlisting}[style=inline, caption={Cargo.toml --- tipo de crate}]
[lib]
crate-type = ["cdylib"]
\end{lstlisting}

O binário final é distribuído via \textbf{GitHub Releases} com CI/CD automático.
O \hayashi{} baixa automaticamente o artefato correto para a arquitetura alvo:

\begin{center}
  \begin{tabular}{@{}ll@{}}
    \toprule
    \textbf{Plataforma}      & \textbf{Arquivo} \\
    \midrule
    Linux x86\_64            & \code{hayplot-x86\_64-unknown-linux-gnu.so} \\
    macOS aarch64            & \code{hayplot-aarch64-apple-darwin.dylib} \\
    Windows x86\_64          & \code{hayplot-x86\_64-pc-windows-msvc.dll} \\
    \bottomrule
  \end{tabular}
\end{center}

% ============================================================
\chapter{Guia de Contribuição}
% ============================================================

\section{Configurando o Ambiente}

\begin{lstlisting}[style=inline, caption={Clonar e compilar localmente}]
git clone https://github.com/sheep-farm/hayplot.git
cd hayplot
cargo build
\end{lstlisting}

Para compilar em modo de lançamento (otimizado):

\begin{lstlisting}[style=inline]
cargo build --release
\end{lstlisting}

O arquivo \code{.so} estará em \code{target/release/libhayplot.so}.

\section{Testando com o Hayashi Local}

Para usar a versão local do plugin (sem instalar pelo CLI):

\begin{lstlisting}[caption={Carregando o .so compilado localmente}]
import("./target/debug/libhayplot.so", as=gg)
\end{lstlisting}

\begin{warnbox}[title={\textbf{Atenção}}]
  O caminho com \code{./target/debug/} é adequado apenas para desenvolvimento.
  Em produção, use sempre \code{import("sheep-farm/hayplot", as=gg)} para
  carregar a versão instalada via \code{hay install}.
\end{warnbox}


\section{Adicionando Novas Geometrias}

O padrão para adicionar uma nova geometria (e.g., \code{geom\_line}) é:

\begin{enumerate}
  \item Criar uma função pública anotada com \code{\#[hayashi\_fn]};
  \item A função recebe \code{mut plot: HashMap<String, HayashiValue>};
  \item Adiciona uma entrada ao vetor \code{"layers"} com a chave \code{"geom"};
  \item Retorna o \code{plot} modificado.
\end{enumerate}

\begin{lstlisting}[caption={Estrutura de uma nova geometria em Rust}]
#[hayashi_fn]
pub fn geom_line(
    mut plot: HashMap<String, HayashiValue>,
    color: String,
    width: f64,
) -> HashMap<String, HayashiValue> {
    if let Some(HayashiValue::List(ref mut layers)) = plot.get_mut("layers") {
        let mut layer = HashMap::new();
        layer.insert("geom".to_string(),  HayashiValue::Str("line".to_string()));
        layer.insert("color".to_string(), HayashiValue::Str(color));
        layer.insert("width".to_string(), HayashiValue::Float(width));
        layers.push(HayashiValue::Dict(layer));
    }
    plot
}
\end{lstlisting}

Em seguida, adicione o tratamento correspondente dentro de \code{render\_svg()}
no loop de layers.

% ============================================================
\chapter{Roadmap}
% ============================================================

\begin{center}
  \rowcolors{2}{haylight}{white}
  \begin{tabularx}{\textwidth}{@{}lXl@{}}
    \toprule
    \textbf{Versão} & \textbf{Feature}                                           & \textbf{Status} \\
    \midrule
    v0.1.x & \code{geom\_point} --- dispers\~{a}o                                  & $\checkmark$ Conclu\'{i}do \\
    v0.1.x & \code{geom\_line} --- s\'{e}ries temporais e tend\^{e}ncia               & $\checkmark$ Conclu\'{i}do \\
    v0.3.0 & \code{geom\_bar} --- barras verticais                              & $\checkmark$ Conclu\'{i}do \\
    v0.3.0 & \code{geom\_histogram} --- histograma                              & $\checkmark$ Conclu\'{i}do \\
    v0.3.0 & \code{geom\_boxplot} --- box-and-whisker                           & $\checkmark$ Conclu\'{i}do \\
    v0.3.0 & \code{geom\_heatmap} --- mapa de calor 2D                           & $\checkmark$ Conclu\'{i}do \\
    v0.4.0 & \code{geom\_area} --- área preenchida sob linha                    & $\checkmark$ Conclu\'{i}do \\
    v0.5.0 & \code{geom\_hline} --- linha horizontal                             & $\checkmark$ Conclu\'{i}do \\
    v0.5.0 & \code{geom\_vline} --- linha vertical                               & $\checkmark$ Conclu\'{i}do \\
    v0.5.0 & \code{geom\_abline} --- linha diagonal                               & $\checkmark$ Conclu\'{i}do \\
    v0.5.0 & \code{geom\_step} --- linhas em degrau                              & $\checkmark$ Conclu\'{i}do \\
    v0.6.0 & Suporte a cores hexadecimais (\code{\#RRGGBB})                     & $\checkmark$ Conclu\'{i}do \\
    v0.6.0 & \code{scale\_x\_log10}, \code{scale\_y\_log10} --- eixos em log    & $\checkmark$ Conclu\'{i}do \\
    v0.7.0 & \code{facet\_wrap}, \code{render\_facets} --- painéis por grupo  & $\checkmark$ Conclu\'{i}do \\
    v0.8.0 & Saída PNG (via feature \code{png\_backend})                        & Planejado \\
    v0.8.0 & \code{geom\_ribbon} --- bandas de intervalo de confiança           & Futuro \\
    v0.8.0 & \code{geom\_smooth} --- regressão LOESS/linear                    & Futuro \\
    v0.9.0 & \code{theme\_*} --- temas de estilo (dark, minimal, classic)       & Futuro \\
    \bottomrule
  \end{tabularx}
\end{center}

% ============================================================
\chapter{Referência Rápida}
% ============================================================

\begin{center}
  \rowcolors{2}{haylight}{white}
  \begin{tabularx}{\textwidth}{@{}L{4.8cm}L{4.0cm}X@{}}
    \toprule
    \textbf{Função} & \textbf{Entrada} & \textbf{Descrição} \\
    \midrule
    \code{hayplot(df, aes)}                        & \code{DataFrame, Dict}              & Inicializa a spec com dados e mapeamento $x/y$ \\[3pt]
    \code{geom\_point(plot, color, size)}          & \code{Dict, String, Float}          & Adiciona camada de dispersão (círculos) \\[3pt]
    \code{geom\_line(plot, color, size)}           & \code{Dict, String, Float}          & Adiciona camada de linha conectando pontos \\[3pt]
    \code{geom\_bar(plot, color, width)}           & \code{Dict, String, Float}          & Adiciona camada de barras verticais \\[3pt]
    \code{geom\_histogram(plot, color, bins)}     & \code{Dict, String, Int}            & Adiciona camada de histograma (distribuição) \\[3pt]
    \code{geom\_boxplot(plot, color, width)}       & \code{Dict, String, Float}          & Adiciona camada de boxplot (quartis, mediana) \\[3pt]
    \code{geom\_heatmap(plot, color, cell\_size)}  & \code{Dict, String, Float}          & Adiciona camada de heatmap (intensidade de cor) \\[3pt]
    \code{geom\_area(plot, color, size)}           & \code{Dict, String, Float}          & Adiciona camada de área preenchida sob linha \\[3pt]
    \code{geom\_hline(plot, color, size, yint)}    & \code{Dict, String, Float, Float}  & Adiciona linha horizontal em yintercept \\[3pt]
    \code{geom\_vline(plot, color, size, xint)}    & \code{Dict, String, Float, Float}  & Adiciona linha vertical em xintercept \\[3pt]
    \code{geom\_abline(plot, color, size, slope, int)} & \code{Dict, String, Float, Float, Float} & Adiciona linha diagonal $y = \text{slope} \cdot x + \text{int}$ \\[3pt]
    \code{geom\_step(plot, color, size, dir)}     & \code{Dict, String, Float, String} & Adiciona linha em degrau (hv/vh) \\[3pt]
    \code{scale\_x\_log10(plot)}                   & \code{Dict}                         & Define eixo X para escala logarítmica base 10 \\[3pt]
    \code{scale\_y\_log10(plot)}                   & \code{Dict}                         & Define eixo Y para escala logarítmica base 10 \\[3pt]
    \code{facet\_wrap(plot, group\_col)}            & \code{Dict, String}                 & Especifica coluna para agrupamento \\[3pt]
    \code{render\_facets(plot)}                   & \code{Dict}                         & Renderiza SVGs separados por grupo \\[3pt]
    \code{labs(plot, title, x, y)}                 & \code{Dict, String, String, String} & Define título e nomes dos eixos \\[3pt]
    \code{render\_svg(plot)}                       & \code{Dict}                         & Renderiza e retorna o SVG como \code{String} \\
    \bottomrule
  \end{tabularx}
\end{center}

\bigskip

\noindent\textbf{Cores disponíveis em todas as geometrias:}

\begin{center}
  \begin{tabular}{@{}llll@{}}
    \code{"red"}     & \code{"blue"}    & \code{"green"}   & \code{"yellow"} \\
    \code{"magenta"} & \code{"cyan"}    & \code{"black"}   & \\
  \end{tabular}
\end{center}

\bigskip

\noindent\textbf{Pipeline mínimo funcional:}

\begin{lstlisting}[caption={Menor script completo}]
import("sheep-farm/hayplot", as=gg)

let df = dataframe({"x": [1.0, 2.0, 3.0], "y": [2.0, 4.0, 6.0]})
let svg = gg::hayplot(df, {"x": "x", "y": "y"})
          |> gg::geom_point("blue", 6.0)
          |> gg::render_svg
write(svg, "out.svg")
\end{lstlisting}

% ── Índice remissivo ─────────────────────────────────────────────────────────
\clearpage
\printindex

\end{document}
